% Problem record op_cbc0bf88b109b122. This TeX file is derived from
% database/problems_json/op_cbc0bf88b109b122.json by scripts/sync-tex.mjs; edit the JSON record
% and rerun the script rather than editing this file.
\section{Ordinary-Petz fidelity remainder for relative-entropy data processing}
\label{sec:op_cbc0bf88b109b122}

\paragraph{Problem.}
Does the ordinary Petz recovery map give a universal fidelity remainder for
monotonicity of quantum relative entropy?  Let
$\mathcal N:\mathcal L(A)\to\mathcal L(B)$ be a finite-dimensional quantum
channel, and let $\rho,\sigma\in\mathcal D(A)$ satisfy
$\operatorname{supp}\rho\subseteq\operatorname{supp}\sigma$.  With inverses
taken on the relevant supports, define the Petz map by
\begin{equation}
  \mathcal P_{\sigma,\mathcal N}(X)
  :=\sigma^{1/2}\mathcal N^{\dagger}\!\left(
      \mathcal N(\sigma)^{-1/2}X
      \mathcal N(\sigma)^{-1/2}
    \right)\sigma^{1/2},
  \label{eq:p69-petz-map}
\end{equation}
where $\mathcal N^{\dagger}$ is the Hilbert--Schmidt adjoint.  Write
$D(\tau\Vert\omega):=\operatorname{Tr}[\tau(\log_2\tau-
\log_2\omega)]$ and use squared fidelity
$F(\tau,\omega):=\lVert\sqrt\tau\sqrt\omega\rVert_1^2$.  The proposed
remainder bound for the map in Eq.~\eqref{eq:p69-petz-map} is
\begin{equation}
  D(\rho\Vert\sigma)
  -D\!\left(\mathcal N(\rho)\middle\Vert\mathcal N(\sigma)\right)
  \stackrel{?}{\geq}
  -\log_2 F\!\left(
    \rho,
    \mathcal P_{\sigma,\mathcal N}(\mathcal N(\rho))
  \right).
  \label{eq:p69-petz-remainder}
\end{equation}
Determine whether Eq.~\eqref{eq:p69-petz-remainder} holds for every such
triple $(\rho,\sigma,\mathcal N)$.

\subsection*{Status}

Solved

\subsection*{Source}

Seshadreesan, Berta, and Wilde explicitly proposed the monotonicity in the
R\'enyi parameter whose endpoint consequence is
Eq.~\eqref{eq:p69-petz-remainder}; Wilde records the same conjectured
ordinary-Petz remainder in Section~12.4
\sourcecite{ref:p69-seshadreesan-berta-wilde}{SBW15},
\sourcecite{ref:p69-wilde}{Wil17}.

\subsection*{Progress}

\begin{itemize}
  \item Bhattacharya disproved Eq.~\eqref{eq:p69-petz-remainder} using the
  diagonal pinching channel $\Phi:M_2(\mathbb C)\to M_2(\mathbb C)$ and the
  density matrices
  \begin{equation}
    A=\begin{pmatrix}
      \tfrac12&\tfrac12\\
      \tfrac12&\tfrac12
    \end{pmatrix},
    \qquad
    B=\begin{pmatrix}
      \tfrac34&-\tfrac14\\
      -\tfrac14&\tfrac14
    \end{pmatrix},
    \qquad
    \Phi(X)=\sum_{j=0}^{1}|j\rangle\!\langle j|X|j\rangle\!\langle j|.
    \label{eq:p69-counterexample}
  \end{equation}
  For $(\rho,\sigma,\mathcal N)=(A,B,\Phi)$ from
  Eq.~\eqref{eq:p69-counterexample}, direct evaluation with
  natural logarithms and root fidelity gives a data-processing loss of
  approximately $1.5191$, while the proposed recovery term is approximately
  $1.5349$.  Since replacing root fidelity by squared fidelity and natural
  logarithms by base-two logarithms rescales both sides consistently, this is
  also a counterexample to the convention in
  Eq.~\eqref{eq:p69-petz-remainder}
  \sourcecite{ref:p69-bhattacharya}{Bha25}.

  \item The counterexample in Eq.~\eqref{eq:p69-counterexample} also rules out
  the full R\'enyi-parameter monotonicity proposed by Seshadreesan, Berta, and
  Wilde, because that monotonicity implies the false endpoint inequality
  Eq.~\eqref{eq:p69-petz-remainder}
  \sourcecite{ref:p69-seshadreesan-berta-wilde}{SBW15},
  \sourcecite{ref:p69-bhattacharya}{Bha25}.
\end{itemize}

\subsection*{References}

\begin{enumerate}[label={},leftmargin=4.8em,labelsep=0.5em]
  \footnotesize
  \item[\textup{[SBW15]}]\label{ref:p69-seshadreesan-berta-wilde}
  K. P. Seshadreesan, M. Berta, and M. M. Wilde,
  ``R\'enyi Squashed Entanglement, Discord, and Relative Entropy
  Differences,'' \emph{Journal of Physics A: Mathematical and Theoretical}
  \textbf{48}, 395303 (2015).
  \href{https://doi.org/10.1088/1751-8113/48/39/395303}{doi:10.1088/1751-8113/48/39/395303};
  \href{https://arxiv.org/abs/1410.1443}{arXiv:1410.1443}.

  \item[\textup{[Wil17]}]\label{ref:p69-wilde}
  M. M. Wilde, \emph{Quantum Information Theory}, 2nd ed., Cambridge
  University Press (2017), Sec.~12.4.
  \href{https://doi.org/10.1017/9781316809976}{doi:10.1017/9781316809976};
  \href{https://arxiv.org/abs/1106.1445}{arXiv:1106.1445}.

  \item[\textup{[Bha25]}]\label{ref:p69-bhattacharya}
  S. Bhattacharya, ``Approximate Recoverability and the Quantum Data
  Processing Inequality,'' arXiv preprint (2023), version~3 revised in 2025.
  \href{https://arxiv.org/abs/2309.02074v3}{arXiv:2309.02074v3}.
\end{enumerate}

\subsection*{Comment}

The answer to Eq.~\eqref{eq:p69-petz-remainder} is negative already in
dimension two.  This general-channel counterexample does not resolve the
more structured conditional-mutual-information inequality in
Problem~\ref{sec:problem-70},
where the channel is a partial trace and the reference state is a marginal
of the state being recovered.

\subsection*{Field}

Quantum Communication

\subsection*{Topic}

Quantum recovery; Quantum relative entropy; Matrix and entropy inequalities

\subsection*{ID}

\texttt{op\_cbc0bf88b109b122}
